\documentclass[a4paper,12pt]{article}
\usepackage[T2A]{fontenc}
\usepackage[utf8]{inputenc}
\usepackage[russian]{babel}
\usepackage{amsmath,amssymb,mathtools}
\renewcommand{\familydefault}{\sfdefault}
\usepackage{geometry}
\geometry{left=24mm,right=24mm,top=25mm,bottom=24mm}
\usepackage{microtype}
\usepackage{icomma}
\usepackage{tikz}
\usetikzlibrary{arrows.meta,positioning,shapes.geometric,calc}
\usepackage{booktabs,tabularx,array}
\usepackage{enumitem}
\usepackage{setspace}
\usepackage{xcolor}
\usepackage{fancyhdr}
\usepackage{hyperref}
\hypersetup{hidelinks,pdftitle={Чек-лист: статистические характеристики и частоты},pdfauthor={Лёвин Артём Александрович}}
\definecolor{accent}{HTML}{008080}
\definecolor{darkaccent}{HTML}{005757}
\definecolor{textgray}{HTML}{30373A}
\definecolor{muted}{HTML}{637073}
\color{textgray}
\setstretch{1.20}
\setlength{\parindent}{0pt}
\setlength{\parskip}{0.55em}
\setlength{\headheight}{14pt}
\pagestyle{fancy}
\fancyhf{}
\fancyhead[L]{\footnotesize\color{muted} Статистические характеристики и частоты}
\fancyhead[R]{\footnotesize\color{muted} Кирилл Зиновьев}
\fancyfoot[C]{\footnotesize\color{muted}\thepage}
\renewcommand{\headrulewidth}{0.35pt}
\renewcommand{\footrulewidth}{0pt}
\renewcommand{\sectionmark}[1]{}
\newcommand{\sect}[1]{\vspace{0.8em}\section*{\color{darkaccent}\large #1}\vspace{-0.6em}\color{accent}\rule{\linewidth}{0.45pt}\color{textgray}\vspace{0.3em}}
\newcommand{\subt}[1]{\vspace{0.45em}{\color{darkaccent}\bfseries #1}\par\vspace{-0.2em}}
\newcommand{\idea}[1]{{\color{darkaccent}\bfseries Ключевая идея.} #1}
\newcommand{\ex}[1]{{\color{darkaccent}\bfseries Пример.} #1}
\newcommand{\task}[1]{\item #1\vspace{0.6em}}
\setlist[enumerate]{leftmargin=*,itemsep=0.55em,topsep=0.45em}
\begin{document}
\thispagestyle{empty}
\begin{center}
\vspace*{33mm}
{\small\color{darkaccent}\MakeUppercase{Математика}}\par
\vspace{16mm}
{\LARGE Чек-лист}\par
\vspace{12mm}
{\Large\color{accent} Статистические характеристики\\[3pt] и частоты}\par
\vspace{12mm}
{\normalsize 7 класс. Повторение и систематизация}\par
\vfill
{\large Кирилл Зиновьев}\par
\vspace{10mm}
{1 ноября 2026 г.}\par
\vspace{21mm}
{\small Лёвин Артём Александрович\\ эксклюзивно для Кирилла Зиновьева}\par
\vspace{18mm}
\begin{tikzpicture}[x=1cm,y=1cm]
\draw[accent,line width=0.7pt] (0,0) .. controls (1.1,.36) and (2.0,-.26) .. (3.2,0)
 .. controls (4.4,.28) and (5.3,-.28) .. (6.5,0)
 .. controls (7.7,.3) and (8.7,-.33) .. (10,0);
\end{tikzpicture}
\end{center}
\clearpage
\sect{1. Числовой ряд и объём выборки}
\textbf{Набор данных} (выборка) --- результаты наблюдений. Одно и то же значение может встретиться несколько раз.

\textbf{Объём выборки} $N$ --- общее число наблюдений, включая повторения.

\ex{В ряду $2,\;3,\;2,\;4,\;3,\;2,\;4,\;3$ содержится $N=8$ чисел. Упорядоченный ряд:}
\[
2,\;2,\;2,\;3,\;3,\;3,\;4,\;4.
\]
\idea{При упорядочивании значения переставляют, сохраняя каждое повторение.}

\sect{2. Среднее арифметическое}
Среднее арифметическое --- сумма всех чисел, делённая на их количество:
\[
\overline{x}=\frac{x_1+x_2+\cdots+x_N}{N}.
\]
\ex{Для чисел $6,\;8,\;10$ среднее равно}
\[
\overline{x}=\frac{6+8+10}{3}=8.
\]
\subt{Если одно значение неизвестно}
Среднее пяти чисел равно $9$, четыре числа --- $4,\;8,\;10,\;12$. Тогда
\[
\frac{4+8+10+12+x}{5}=9
\quad\Longrightarrow\quad 34+x=45
\quad\Longrightarrow\quad x=11.
\]
Проверка: сумма всех пяти чисел равна $45$, а $45:5=9$.

\clearpage
\sect{3. Медиана}
\textbf{Медиана} --- центральное значение упорядоченного ряда. Если наблюдений чётное количество, берут среднее арифметическое двух центральных чисел.

\ex{Нечётное число элементов:}
\[
1,\;\;\; \boxed{4},\;\;\;7
\qquad\Longrightarrow\qquad \text{медиана}=4.
\]
\ex{Чётное число элементов:}
\[
1,\;\;\;\underbrace{4,\;6}_{\text{центральные}},\;\;\;9
\qquad\Longrightarrow\qquad \text{медиана}=\frac{4+6}{2}=5.
\]
\idea{Сначала упорядочьте все данные. Медиану нельзя надёжно находить по неупорядоченному ряду.}

\sect{4. Размах}
\textbf{Размах} --- разность наибольшего и наименьшего значения:
\[
R=x_{\max}-x_{\min}.
\]
\ex{В наборе $-4,\;2,\;6,\;x$ число $x$ --- наибольшее, а размах равен $15$. Тогда}
\[
x-(-4)=15\quad\Longrightarrow\quad x+4=15\quad\Longrightarrow\quad x=11.
\]
Проверка: $11-(-4)=15$; действительно, $11>6$.

\clearpage
\sect{5. Как влияют необычно большие значения}
Рассмотрим числа $3,\;4,\;5,\;6$ и добавим число $37$.
\begin{center}
\renewcommand{\arraystretch}{1.28}
\begin{tabularx}{\linewidth}{@{}Xcc@{}}
\toprule
& До добавления & После добавления \\
\midrule
Среднее арифметическое & $4{,}5$ & $11$ \\
Медиана & $4{,}5$ & $5$ \\
\bottomrule
\end{tabularx}
\end{center}
\textbf{Вывод:} большое значение изменило среднее на $6{,}5$, медиану --- на $0{,}5$. Среднее арифметическое часто сильнее реагирует на крайние значения.

\textbf{Важно:} медиана --- центральное значение, поэтому она \emph{не обязана} совпадать с числом, встречающимся чаще всего.

\sect{6. Абсолютная частота}
\textbf{Абсолютная частота} $n_i$ показывает, сколько раз значение $x_i$ встретилось в ряду. Индекс $i$ --- номер значения в таблице.

\ex{Для набора $2,\;3,\;2,\;4,\;3,\;2,\;4,\;3$ получаем:}
\begin{center}
\renewcommand{\arraystretch}{1.25}
\begin{tabular}{@{}lccc@{}}
\toprule
Значение $x_i$ & $2$ & $3$ & $4$ \\
\midrule
Частота $n_i$ & $3$ & $3$ & $2$ \\
\bottomrule
\end{tabular}
\end{center}
\[
N=n_1+n_2+\cdots+n_k=3+3+2=8.
\]
\idea{Число различных значений здесь равно $3$, объём выборки равен $8$. Не путайте их.}

\clearpage
\sect{7. Относительная частота}
\textbf{Относительная частота} $\omega_i$ --- доля всех наблюдений, приходящаяся на значение $x_i$:
\[
\boxed{\omega_i=\frac{n_i}{N}}.
\]
В процентах:
\[
\boxed{p_i=\frac{n_i}{N}\cdot100\%}.
\]
\textbf{Величины:} $n_i$ --- число повторений, $N$ --- объём выборки, $p_i$ --- процент наблюдений с данным значением.

\subt{Полная таблица распределения}
Для предыдущего ряда $N=8$:
\begin{center}
\renewcommand{\arraystretch}{1.35}
\begin{tabularx}{\linewidth}{@{}Xccc@{}}
\toprule
Значение $x_i$ & $2$ & $3$ & $4$ \\
\midrule
Абсолютная частота $n_i$ & $3$ & $3$ & $2$ \\
Относительная частота $\omega_i$ & $\frac{3}{8}$ & $\frac{3}{8}$ & $\frac{2}{8}$ \\
Относительная частота, \% & $37{,}5$ & $37{,}5$ & $25$ \\
\bottomrule
\end{tabularx}
\end{center}
\[
\frac{3}{8}=0{,}375=37{,}5\%,\qquad \frac{2}{8}=0{,}25=25\%.
\]
\textbf{Две обязательные проверки:}
\[
\sum_i n_i=N,\qquad \sum_i\omega_i=1
\qquad \text{или}\qquad \sum_i p_i=100\%.
\]
Для примера: $3+3+2=8$ и $37{,}5+37{,}5+25=100$.

\subt{Типичные ошибки}
\begin{enumerate}[label=\arabic*. ,itemsep=0.1em]
\item В знаменателе берут число \emph{различных} значений вместо общего числа наблюдений.
\item Теряют повторы при переписывании ряда.
\item Пишут $3{,}75\%$ вместо $37{,}5\%$: при переводе $0{,}375$ в проценты нужно умножить на $100$.
\item Забывают убедиться, что сумма всех абсолютных частот равна $N$.
\end{enumerate}

\clearpage
\begin{center}
{\Large\color{darkaccent} Алгоритм: составление таблицы частот}\par
\vspace{11mm}
\end{center}
\begin{center}
\begin{tikzpicture}[
 node distance=11mm,
 start/.style={ellipse,draw=darkaccent,very thick,align=center,minimum width=36mm,minimum height=10mm,font=\small},
 action/.style={rectangle,rounded corners=2pt,draw=accent,thick,align=center,text width=85mm,minimum height=13mm,font=\small,inner sep=4pt},
 check/.style={diamond,aspect=2.0,draw=darkaccent,thick,align=center,text width=47mm,inner sep=1pt,font=\small},
 arr/.style={-{Stealth[length=2.5mm]},thick,draw=darkaccent}
 ]
 \node[start] (begin) {Исходный ряд};
 \node[action,below=of begin] (sort) {Упорядочить значения и подсчитать\\общее число наблюдений $N$};
 \node[action,below=of sort] (freq) {Перечислить разные значения $x_i$\\и подсчитать их частоты $n_i$};
 \node[action,below=of freq] (rel) {Найти $\omega_i=n_i/N$\\и $p_i=\omega_i\cdot100\%$};
 \node[check,below=14mm of rel] (test) {Сумма $n_i=N$\\и сумма $p_i=100\%$?};
 \node[start,below=16mm of test] (end) {Таблица готова};
 \draw[arr] (begin)--(sort);
 \draw[arr] (sort)--(freq);
 \draw[arr] (freq)--(rel);
 \draw[arr] (rel)--(test);
 \draw[arr] (test)--node[right,font=\footnotesize]{да}(end);
 \coordinate (goLeft) at ($(test.west)+(-22mm,0)$);
 \coordinate (toFreq) at ($(freq.west)+(-22mm,0)$);
 \draw[arr] (test.west)--node[above,font=\footnotesize]{нет}(goLeft) |- (toFreq) -- (freq.west);
\end{tikzpicture}
\end{center}
\clearpage
\sect{8. Короткая памятка перед решением}
\begin{tabularx}{\linewidth}{@{}>{\color{darkaccent}\bfseries}p{0.35\linewidth}X@{}}
\toprule
Что нужно найти & Что сделать \\
\midrule
Объём выборки & Посчитать все наблюдения, включая повторения. \\
Среднее & Сложить все числа и разделить на их количество. \\
Медиану & Упорядочить ряд; взять центр или среднее двух центральных. \\
Размах & Вычесть минимум из максимума. \\
Абсолютную частоту & Подсчитать повторения нужного значения. \\
Относительную частоту & Разделить частоту на объём выборки; при необходимости перевести в проценты. \\
Проверку таблицы & Сложить частоты: $N$; сложить проценты: $100\%$. \\
\bottomrule
\end{tabularx}

\vspace{0.6em}
\textbf{Когда абсолютные частоты уже даны в таблице,} среднее можно считать без выписывания всего ряда:
\[
\overline{x}=\frac{x_1n_1+x_2n_2+\cdots+x_kn_k}{N}.
\]
Это тот же способ нахождения среднего: каждое значение прибавляется столько раз, сколько оно встречается.

\ex{Значения $2,\;3,\;4$ встречаются $3,\;3,\;2$ раза. Тогда}
\[
\overline{x}=\frac{2\cdot3+3\cdot3+4\cdot2}{8}
=\frac{23}{8}=2{,}875.
\]

\clearpage
\sect{Тренировочный блок. Путь А}
Во всех заданиях показывайте вычисления. При составлении таблиц записывайте значения по возрастанию.
\subt{Базовый уровень}
\begin{enumerate}[label=\textbf{\arabic*.}]
\task{Вычислите среднее арифметическое чисел $7,\;11,\;12$.}
\task{Найдите медиану и размах ряда $4,\;2,\;9,\;2,\;8,\;6$.}
\task{Среднее пяти чисел равно $8$. Известны четыре числа: $4,\;6,\;8,\;12$. Найдите пятое число.}
\subt{Средний уровень}
\task{Числа $2,\;5,\;8,\;x,\;13,\;16$ записаны по возрастанию. Их медиана равна $10$. Найдите $x$.}
\task{Дан ряд $3,\;4,\;5,\;6$. К нему добавили число $37$. Вычислите среднее арифметическое и медиану \emph{до и после} добавления. Какая характеристика изменилась сильнее?}
\task{Составьте таблицу абсолютных частот для ряда\\ $2,\;4,\;2,\;3,\;4,\;2,\;3,\;3,\;3,\;4,\;4,\;4$. Укажите объём выборки.}
\task{Даны абсолютные частоты значений $1,\;2,\;3$: соответственно $2,\;3,\;5$. Найдите объём выборки и относительные частоты в процентах.}
\subt{Уровень выше среднего}
\task{Для ряда $2,\;3,\;2,\;4,\;3,\;2,\;4,\;3$ составьте полную таблицу: значение, абсолютная частота, относительная частота (дробь и проценты). Проверьте обе суммы.}
\task{Объём выборки $20$. Значение $1$ встретилось $6$ раз, значение $3$ --- $5$ раз, значение $2$ составляет $20\%$ выборки. Все остальные наблюдения равны $4$. Найдите абсолютные и относительные частоты всех четырёх значений.}
\task{Дана таблица: значение $2$ встретилось $2$ раза, значение $4$ --- $3$ раза, значение $8$ --- $5$ раз. Вычислите объём выборки, относительные частоты (в процентах), среднее арифметическое, медиану и размах.}
\end{enumerate}

\clearpage
\sect{Ответы к тренировочному блоку}
\renewcommand{\arraystretch}{1.42}
\begin{tabularx}{\linewidth}{@{}>{\bfseries}p{14mm}X@{}}
\toprule
№ & Ответ \\
\midrule
1 & $10$. \\
2 & Медиана $5$; размах $7$. \\
3 & $10$. \\
4 & $x=12$. \\
5 & До: среднее $4{,}5$, медиана $4{,}5$; после: среднее $11$, медиана $5$. Сильнее изменилось среднее. \\
6 & $N=12$; частоты: $n(2)=3$, $n(3)=4$, $n(4)=5$. \\
7 & $N=10$; $20\%$, $30\%$, $50\%$. \\
8 & $N=8$; для $2$: $3$, $\frac38$, $37{,}5\%$; для $3$: $3$, $\frac38$, $37{,}5\%$; для $4$: $2$, $\frac28$, $25\%$. Суммы: $8$ и $100\%$. \\
9 & Частоты: $6,4,5,5$; проценты: $30\%,20\%,25\%,25\%$. \\
10 & $N=10$; проценты: $20\%,30\%,50\%$; среднее $5{,}6$; медиана $6$; размах $6$. \\
\bottomrule
\end{tabularx}
\vfill
\begin{center}
\color{muted}\footnotesize Лёвин Артём Александрович\\ эксклюзивно для Кирилла Зиновьева
\end{center}
\end{document}
